\section{The difference in Morse counts}\label{sec:counts}

The four fixed components of the involution on $X\times Y$ are copies
of $X\times\CP^1$.  We now compare their combined Morse count with the
stable Morse number of the ambient manifold.  Set
\begin{equation}\label{eq:counts-products}
 C=X\times\CP^1,\qquad M=X\times Y.
\end{equation}
The calculation rests on the fact that torsion at distinct primes can
share a generator.  For nonnegative integers $a,q$,
\begin{equation}\label{eq:counts-generators}
 \text{the minimum number of generators of }
 (\Z/2)^a\oplus(\Z/3)^q
 \text{ is }\max(a,q).
\end{equation}
Reduction modulo $2$ and modulo $3$ gives the lower bound; pairing
cyclic summands of different orders into copies of $\Z/6$ gives the
upper bound.

\begin{proposition}\label{prop:gap}
For the manifolds in \eqref{eq:counts-products},
\begin{equation}\label{eq:counts-values}
 \lambda(C)=114N-518,\qquad
 \lambda(M)=456N-2056.
\end{equation}
In particular,
\begin{equation}\label{eq:counts-gap}
 \lambda(M)-4\lambda(C)=16.
\end{equation}
\end{proposition}

\begin{proof}
Recall from Section~\ref{sec:geometry} that $L=N-3$ and that, in each of the two lists of degrees
$1+2j$ and $2+2j$, the torsion in $H_*(X;\Z)$ has multiplicities
\[
 a_j=2,\qquad
 (q_1,\ldots,q_L)=(1,3,4,\ldots,4,3,1),
 \qquad 1\leq j\leq L,
\]
at the primes $2$ and $3$, respectively.  Extend both sequences by
zero to all $j\in\Z$.  Put $c_j=\max(a_j,q_j)$ and
\begin{equation}\label{eq:counts-S}
 S=\sum_jc_j=2+3+4(L-4)+3+2=4L-6=4N-18.
\end{equation}

The homology of $\CP^1$ is free with ranks $(1,1)$ in degrees
$(0,2)$, and that of $Y$ is free with ranks $(1,6,1)$ in degrees
$(0,2,4)$.  Thus the integral K\"unneth formula has no torsion-product
terms in either product.  It replaces the profiles $a,q$ by their
convolutions with
\[
 e_C=(1,1),\qquad e_M=(1,6,1),
\]
where these sequences start at index $0$ and
$(e*a)_j=\sum_ke_ka_{j-k}$.  By
\eqref{eq:counts-generators}, the sum of torsion generator numbers
in either one of the two degree lists is therefore
\[
 \sum_j\max\bigl((e*a)_j,(e*q)_j\bigr).
\]

We compute how much this sum falls short of
$\sum_j(e*c)_j=(\sum_ke_k)S$.  Define the nonnegative sequences
\[
 u_j=\max(a_j-q_j,0),\qquad
 v_j=\max(q_j-a_j,0).
\]
Since $c=a+v=q+u$, the loss at index $j$ is exactly
\begin{equation}\label{eq:counts-loss}
 (e*c)_j-\max\bigl((e*a)_j,(e*q)_j\bigr)
 =\min\bigl((e*u)_j,(e*v)_j\bigr).
\end{equation}
Here $u_1=u_L=1$ and $u_j=0$ elsewhere, whereas
$v_2=v_{L-1}=1$, $v_j=2$ for $3\leq j\leq L-2$, and $v_j=0$
elsewhere.  Consequently the loss can occur only near the two ends
of the profiles.  Table~\ref{tab:counts-loss} gives every index at
which $(e*u)_j$ is nonzero.  The ends do not overlap, since $L\geq7$.

\begin{table}[htbp]
\centering
\begin{tabular}{@{}ccrrr@{}}
\toprule
Kernel $e$ & Indices $j$ & $(e*u)_j$ & $(e*v)_j$ & Loss \\
\midrule
$(1,1)$ & $1,\ L+1$ & $1$ & $0$ & $0$ \\
         & $2,\ L$   & $1$ & $1$ & $1$ \\
\midrule
$(1,6,1)$ & $1,\ L+2$ & $1$ & $0$ & $0$ \\
           & $2,\ L+1$ & $6$ & $1$ & $1$ \\
           & $3,\ L$   & $1$ & $8$ & $1$ \\
\bottomrule
\end{tabular}
\caption{The loss in the torsion generator count under convolution.
Both indices in each pair have the displayed values; every unlisted
index has zero loss.}
\label{tab:counts-loss}
\end{table}

The total loss is $2$ for $e_C$ and $4$ for $e_M$.  All degree shifts
are even, so the two parity lists remain separate and give identical
contributions.  Hence
\begin{equation}\label{eq:counts-total-torsion}
 \sum_i t_i(C)=2(2S-2),\qquad
 \sum_i t_i(M)=2(8S-4).
\end{equation}
The free ranks satisfy $b(C)=2b(X)$ and $b(M)=8b(X)$.  It follows
that
\begin{align*}
 \lambda(C)&=2b(X)+8S-8,\\
 \lambda(M)&=8b(X)+32S-16.
\end{align*}
Subtracting four times the first identity from the second proves
\eqref{eq:counts-gap}.  Finally, substituting $b(X)=41N-183$ and
$S=4N-18$ gives \eqref{eq:counts-values}.
\end{proof}

It remains to realize $4\lambda(C)$ as a Hamiltonian fixed-point count on $M$.
